\def\probtitle{덜 개성있는 구간}
\def\probno{F} % 문제 번호

\begin{problem}{\probtitle}{표준 입력(stdin)}{표준 출력(stdout)}{1 초}{1024 megabytes}

다음 조건을 만족하는 길이가 $N$인 수열 $A = [A_1, A_2, \cdots, A_N]$를 아무거나 하나 찾아라.

\begin{itemize}[topsep=0pt,noitemsep]
    \item 수열 $A$의 모든 원소는 서로 다르다.
    \item $1 \le L \le R \le N$을 만족하는 모든 정수 쌍 $(L, R)$에 대하여, $\sum\limits_{i=L}^R A_i$가 가질 수 있는 서로 다른 값의 개수는 정확히 $N$개이다.
\end{itemize}

주어진 범위 내에서 조건을 만족하는 수열은 항상 존재한다.

\InputFile

첫째 줄에 수열 $A$의 길이를 나타내는 정수 $N$이 주어진다.

\OutputFile

첫째 줄에 조건을 만족하는 수열 $A$의 원소들을 공백으로 구분하여 출력한다.

위 조건을 만족하는 수열이 여러 개라면 그중 아무거나 출력한다.

\Constraints

\begin{itemize}[topsep=0pt,noitemsep]
    \item 주어지는 모든 수는 정수이다.
    \item $1 \le N \le 5\, 000$
    \item $1 \le i \le N$인 모든 $i$에 대해 $-10^{9} \le A_i \le 10^{9}$이다.
\end{itemize}

\Example

\begin{example}
    \exmpfile{./example/01.in.txt}{./example/01.out.txt}%
    \exmpfile{./example/02.in.txt}{./example/02.out.txt}%
\end{example}

\Notes

수열 $A = [-2, 0, 2]$에서 모든 연속 구간의 합을 구하면 다음과 같다.

\begin{itemize}[topsep=0pt,noitemsep]
    \item $-2$: $(1, 1)$, $(1, 2)$
    \item $0$: $(1, 3)$, $(2, 2)$
    \item $2$: $(2, 3)$, $(3, 3)$
\end{itemize}

\end{problem}